\documentclass{article}%
\usepackage[T1]{fontenc}%
\usepackage[utf8]{inputenc}%
\usepackage{lmodern}%
\usepackage{textcomp}%
\usepackage{lastpage}%
\usepackage{geometry}%
\geometry{margin=1in,headheight=14pt,headsep=25pt}%
\usepackage{hyperref}%
\usepackage{enumitem}%
\usepackage{fancyhdr}%
\usepackage{xcolor}%
\usepackage{xurl}%
\usepackage{breakurl}%
\usepackage{amsmath}%
\usepackage{amssymb}%
\usepackage{mathtools}%
\usepackage{amsthm}%
\usepackage{thmtools}%
%

            \hypersetup{
                colorlinks=true,
                linkcolor=blue,
                filecolor=magenta,
                urlcolor=blue,
                breaklinks=true,
                bookmarks=true,
                pdfborder={0 0 0}
            }
            \urlstyle{same}
            \Urlmuskip=0mu plus 1mu  % URL breaking in nicer spots
            
            % Math configuration
            \everymath{\displaystyle}
            \setlength{\jot}{10pt}  % Increase spacing between equations
            
            \pagestyle{fancy}
            \fancyhf{}
            \rhead{Generated on \today}
            \lhead{Course Summary}
            \cfoot{\thepage}
            
            % Better Unicode support
            \usepackage[utf8]{inputenc}
            \usepackage{newunicodechar}
            
            % Configure itemize settings globally
            \setlist[itemize]{leftmargin=*,nosep}
        %
\title{Artifical Intelligence Basics Course Summary}%
\author{Generated by Scribe.AI}%
\date{\today}%
%
\begin{document}%
\normalsize%
\maketitle%
\section{Key Terms}%
\label{sec:KeyTerms}%
\subsection{Artificial Neural Networks}%
\label{subsec:ArtificialNeuralNetworks}%
\begin{itemize}[leftmargin=*]%
\item \href{\url{https://purdue.brightspace.com/d2l/common/assets/pdfjs-d2l-dist/1.0.14-legacy/web/viewer.html?file=\%252Fcontent\%252Fenforced\%252F1095467-wl.202510.CS.24300.001\%252F02L2-Intro.pdf\%253Fou\%253D1095467\&lang=en-us\&container=d2l-fileviewer-rendered-pdf\&fullscreen=d2l-fileviewer-rendered-pdf-dialog\&height=667}#page=6}{\textbf{Turing Test}}: An operational test for intelligent behavior, the Imitation Game, where an interrogator tries to distinguish between a human and a machine based on their responses.%
\item \href{\url{https://purdue.brightspace.com/d2l/common/assets/pdfjs-d2l-dist/1.0.14-legacy/web/viewer.html?file=\%252Fcontent\%252Fenforced\%252F1095467-wl.202510.CS.24300.001\%252F02L2-Intro.pdf\%253Fou\%253D1095467\&lang=en-us\&container=d2l-fileviewer-rendered-pdf\&fullscreen=d2l-fileviewer-rendered-pdf-dialog\&height=667}#page=45}{\textbf{Perceptron}}: A fundamental unit in neural networks that takes multiple inputs, applies weights to them, sums the weighted inputs, and produces an output based on a threshold.%
\end{itemize}

%
\subsection{Machine Learning Algorithms}%
\label{subsec:MachineLearningAlgorithms}%
\begin{itemize}[leftmargin=*]%
\item \href{\url{https://purdue.brightspace.com/d2l/common/assets/pdfjs-d2l-dist/1.0.14-legacy/web/viewer.html?file=\%252Fcontent\%252Fenforced\%252F1095467-wl.202510.CS.24300.001\%252F02L2-Intro.pdf\%253Fou\%253D1095467\&lang=en-us\&container=d2l-fileviewer-rendered-pdf\&fullscreen=d2l-fileviewer-rendered-pdf-dialog\&height=667}#page=46}{\textbf{Reinforcement Learning}}: A type of machine learning where an agent learns to make decisions by interacting with an environment and receiving rewards or penalties.%
\end{itemize}

%
\subsection{Search Algorithms}%
\label{subsec:SearchAlgorithms}%
\begin{itemize}[leftmargin=*]%
\item \href{\url{https://purdue.brightspace.com/d2l/common/assets/pdfjs-d2l-dist/1.0.14-legacy/web/viewer.html?file=\%252Fcontent\%252Fenforced\%252F1095467-wl.202510.CS.24300.001\%252F02L2-Intro.pdf\%253Fou\%253D1095467\&lang=en-us\&container=d2l-fileviewer-rendered-pdf\&fullscreen=d2l-fileviewer-rendered-pdf-dialog\&height=667}#page=22}{\textbf{Minimax Search}}: A search algorithm for two-player zero-sum games that aims to find the best move for the current player, assuming the opponent also plays optimally.%
\item \href{\url{https://purdue.brightspace.com/d2l/common/assets/pdfjs-d2l-dist/1.0.14-legacy/web/viewer.html?file=\%252Fcontent\%252Fenforced\%252F1095467-wl.202510.CS.24300.001\%252F02L2-Intro.pdf\%253Fou\%253D1095467\&lang=en-us\&container=d2l-fileviewer-rendered-pdf\&fullscreen=d2l-fileviewer-rendered-pdf-dialog\&height=667}#page=23}{\textbf{Alpha-Beta Pruning}}: A technique to reduce the search space in Minimax search by avoiding the evaluation of branches that cannot affect the optimal decision.%
\end{itemize}

%
\subsection{Linear Algebra Fundamentals}%
\label{subsec:LinearAlgebraFundamentals}%
\begin{itemize}[leftmargin=*]%
\item \href{\url{https://purdue.brightspace.com/d2l/common/assets/pdfjs-d2l-dist/1.0.14-legacy/web/viewer.html?file=\%252Fcontent\%252Fenforced\%252F1095467-wl.202510.CS.24300.001\%252F04L1-Linear\%20Algebra\%20Primer.pdf\%253Fou\%253D1095467\&lang=en-us\&container=d2l-fileviewer-rendered-pdf\&fullscreen=d2l-fileviewer-rendered-pdf-dialog\&height=667}#page=4}{\textbf{Matrix}}: A rectangular array of numbers, symbols, or expressions arranged in rows and columns.%
\item \href{\url{https://purdue.brightspace.com/d2l/common/assets/pdfjs-d2l-dist/1.0.14-legacy/web/viewer.html?file=\%252Fcontent\%252Fenforced\%252F1095467-wl.202510.CS.24300.001\%252F04L1-Linear\%20Algebra\%20Primer.pdf\%253Fou\%253D1095467\&lang=en-us\&container=d2l-fileviewer-rendered-pdf\&fullscreen=d2l-fileviewer-rendered-pdf-dialog\&height=667}#page=5}{\textbf{Linear Transform}}: A function that maps vectors from one vector space to another, satisfying $T(u + v) = T(u) + T(v)$ and $T(cu) = cT(u)$ for all vectors $u$, $v$ and scalars $c$.%
\item \href{\url{https://purdue.brightspace.com/d2l/common/assets/pdfjs-d2l-dist/1.0.14-legacy/web/viewer.html?file=\%252Fcontent\%252Fenforced\%252F1095467-wl.202510.CS.24300.001\%252F04L1-Linear\%20Algebra\%20Primer.pdf\%253Fou\%253D1095467\&lang=en-us\&container=d2l-fileviewer-rendered-pdf\&fullscreen=d2l-fileviewer-rendered-pdf-dialog\&height=667}#page=14}{\textbf{Matrix Multiplication}}: An operation that combines two matrices to produce a third matrix, where the element at row $i$ and column $j$ of the resulting matrix is the dot product of row $i$ of the first matrix and column $j$ of the second matrix. The inner dimensions must match.%
\item \href{\url{https://purdue.brightspace.com/d2l/common/assets/pdfjs-d2l-dist/1.0.14-legacy/web/viewer.html?file=\%252Fcontent\%252Fenforced\%252F1095467-wl.202510.CS.24300.001\%252F04L1-Linear\%20Algebra\%20Primer.pdf\%253Fou\%253D1095467\&lang=en-us\&container=d2l-fileviewer-rendered-pdf\&fullscreen=d2l-fileviewer-rendered-pdf-dialog\&height=667}#page=10}{\textbf{Matrix Inverse}}: For a square matrix $A$, its inverse $A^{-1}$ is a matrix such that $AA^{-1} = A^{-1}A = I$, where $I$ is the identity matrix.%
\item \href{\url{https://purdue.brightspace.com/d2l/common/assets/pdfjs-d2l-dist/1.0.14-legacy/web/viewer.html?file=\%252Fcontent\%252Fenforced\%252F1095467-wl.202510.CS.24300.001\%252F04L1-Linear\%20Algebra\%20Primer.pdf\%253Fou\%253D1095467\&lang=en-us\&container=d2l-fileviewer-rendered-pdf\&fullscreen=d2l-fileviewer-rendered-pdf-dialog\&height=667}#page=16}{\textbf{Matrix Transposition}}: An operation that swaps the rows and columns of a matrix. The transpose of matrix $A$ is denoted as $A^T$.%
\item \href{\url{https://purdue.brightspace.com/d2l/common/assets/pdfjs-d2l-dist/1.0.14-legacy/web/viewer.html?file=\%252Fcontent\%252Fenforced\%252F1095467-wl.202510.CS.24300.001\%252F04L1-Linear\%20Algebra\%20Primer.pdf\%253Fou\%253D1095467\&lang=en-us\&container=d2l-fileviewer-rendered-pdf\&fullscreen=d2l-fileviewer-rendered-pdf-dialog\&height=667}#page=21}{\textbf{Derivative}}: A measure of how a function changes in response to a small change in its input. It is the slope of the tangent line to the function at a given point.%
\item \href{\url{https://purdue.brightspace.com/d2l/common/assets/pdfjs-d2l-dist/1.0.14-legacy/web/viewer.html?file=\%252Fcontent\%252Fenforced\%252F1095467-wl.202510.CS.24300.001\%252F04L1-Linear\%20Algebra\%20Primer.pdf\%253Fou\%253D1095467\&lang=en-us\&container=d2l-fileviewer-rendered-pdf\&fullscreen=d2l-fileviewer-rendered-pdf-dialog\&height=667}#page=19}{\textbf{Partial Derivative}}: The derivative of a function with respect to one of its variables, holding all other variables constant.%
\item \href{\url{https://purdue.brightspace.com/d2l/common/assets/pdfjs-d2l-dist/1.0.14-legacy/web/viewer.html?file=\%252Fcontent\%252Fenforced\%252F1095467-wl.202510.CS.24300.001\%252F04L1-Linear\%20Algebra\%20Primer.pdf\%253Fou\%253D1095467\&lang=en-us\&container=d2l-fileviewer-rendered-pdf\&fullscreen=d2l-fileviewer-rendered-pdf-dialog\&height=667}#page=21}{\textbf{Sigmoid Function}}: A mathematical function that maps any input value to an output value between 0 and 1. Often used as an activation function in neural networks.%
\end{itemize}

%
\subsection{Game Playing AI}%
\label{subsec:GamePlayingAI}%
\begin{itemize}[leftmargin=*]%
\item \href{\url{https://purdue.brightspace.com/d2l/common/assets/pdfjs-d2l-dist/1.0.14-legacy/web/viewer.html?file=\%252Fcontent\%252Fenforced\%252F1095467-wl.202510.CS.24300.001\%252F02L2-Intro.pdf\%253Fou\%253D1095467\&lang=en-us\&container=d2l-fileviewer-rendered-pdf\&fullscreen=d2l-fileviewer-rendered-pdf-dialog\&height=667}#page=11}{\textbf{Game Tree}}: A tree-like structure representing all possible sequences of moves in a game, with leaf nodes representing outcomes.%
\end{itemize}

%
\subsection{Inductive Reasoning}%
\label{subsec:InductiveReasoning}%
\begin{itemize}[leftmargin=*]%
\item \href{\url{https://purdue.brightspace.com/d2l/common/assets/pdfjs-d2l-dist/1.0.14-legacy/web/viewer.html?file=\%252Fcontent\%252Fenforced\%252F1095467-wl.202510.CS.24300.001\%252F05L2-ML\%20Basics.pdf\%253Fou\%253D1095467\&lang=en-us\&container=d2l-fileviewer-rendered-pdf\&fullscreen=d2l-fileviewer-rendered-pdf-dialog\&height=667}#page=3}{\textbf{Inductive Reasoning}}: The process of deriving general principles from specific observations.%
\item \href{\url{https://purdue.brightspace.com/d2l/common/assets/pdfjs-d2l-dist/1.0.14-legacy/web/viewer.html?file=\%252Fcontent\%252Fenforced\%252F1095467-wl.202510.CS.24300.001\%252F05L2-ML\%20Basics.pdf\%253Fou\%253D1095467\&lang=en-us\&container=d2l-fileviewer-rendered-pdf\&fullscreen=d2l-fileviewer-rendered-pdf-dialog\&height=667}#page=6}{\textbf{Linear Regression}}: A machine learning algorithm used to model the relationship between a dependent variable and one or more independent variables by fitting a linear equation to observed data.%
\item \href{\url{https://purdue.brightspace.com/d2l/common/assets/pdfjs-d2l-dist/1.0.14-legacy/web/viewer.html?file=\%252Fcontent\%252Fenforced\%252F1095467-wl.202510.CS.24300.001\%252F05L2-ML\%20Basics.pdf\%253Fou\%253D1095467\&lang=en-us\&container=d2l-fileviewer-rendered-pdf\&fullscreen=d2l-fileviewer-rendered-pdf-dialog\&height=667}#page=11}{\textbf{Multiple Linear Regression}}: An extension of linear regression that models the relationship between a dependent variable and two or more independent variables.%
\item \href{\url{https://purdue.brightspace.com/d2l/common/assets/pdfjs-d2l-dist/1.0.14-legacy/web/viewer.html?file=\%252Fcontent\%252Fenforced\%252F1095467-wl.202510.CS.24300.001\%252F05L2-ML\%20Basics.pdf\%253Fou\%253D1095467\&lang=en-us\&container=d2l-fileviewer-rendered-pdf\&fullscreen=d2l-fileviewer-rendered-pdf-dialog\&height=667}#page=14}{\textbf{Overfitting}}: A phenomenon in machine learning where a model learns the training data too well, resulting in poor generalization to unseen data.%
\item \href{\url{https://purdue.brightspace.com/d2l/common/assets/pdfjs-d2l-dist/1.0.14-legacy/web/viewer.html?file=\%252Fcontent\%252Fenforced\%252F1095467-wl.202510.CS.24300.001\%252F05L2-ML\%20Basics.pdf\%253Fou\%253D1095467\&lang=en-us\&container=d2l-fileviewer-rendered-pdf\&fullscreen=d2l-fileviewer-rendered-pdf-dialog\&height=667}#page=20}{\textbf{Binary Classification}}: A type of machine learning problem where the goal is to classify data into one of two categories.%
\item \href{\url{https://purdue.brightspace.com/d2l/common/assets/pdfjs-d2l-dist/1.0.14-legacy/web/viewer.html?file=\%252Fcontent\%252Fenforced\%252F1095467-wl.202510.CS.24300.001\%252F05L2-ML\%20Basics.pdf\%253Fou\%253D1095467\&lang=en-us\&container=d2l-fileviewer-rendered-pdf\&fullscreen=d2l-fileviewer-rendered-pdf-dialog\&height=667}#page=24}{\textbf{Decision Boundary}}: A line or hyperplane that separates data points into different classes in a classification problem.%
\item \href{\url{https://purdue.brightspace.com/d2l/common/assets/pdfjs-d2l-dist/1.0.14-legacy/web/viewer.html?file=\%252Fcontent\%252Fenforced\%252F1095467-wl.202510.CS.24300.001\%252F05L2-ML\%20Basics.pdf\%253Fou\%253D1095467\&lang=en-us\&container=d2l-fileviewer-rendered-pdf\&fullscreen=d2l-fileviewer-rendered-pdf-dialog\&height=667}#page=25}{\textbf{0-1 Loss}}: A loss function in binary classification that assigns a loss of 1 if the prediction is incorrect and 0 if it is correct.%
\item \href{\url{https://purdue.brightspace.com/d2l/common/assets/pdfjs-d2l-dist/1.0.14-legacy/web/viewer.html?file=\%252Fcontent\%252Fenforced\%252F1095467-wl.202510.CS.24300.001\%252F05L2-ML\%20Basics.pdf\%253Fou\%253D1095467\&lang=en-us\&container=d2l-fileviewer-rendered-pdf\&fullscreen=d2l-fileviewer-rendered-pdf-dialog\&height=667}#page=27}{\textbf{Perceptron Loss}}: A loss function used in binary classification that is differentiable and encourages correct classification.%
\item \href{\url{https://purdue.brightspace.com/d2l/common/assets/pdfjs-d2l-dist/1.0.14-legacy/web/viewer.html?file=\%252Fcontent\%252Fenforced\%252F1095467-wl.202510.CS.24300.001\%252F05L2-ML\%20Basics.pdf\%253Fou\%253D1095467\&lang=en-us\&container=d2l-fileviewer-rendered-pdf\&fullscreen=d2l-fileviewer-rendered-pdf-dialog\&height=667}#page=27}{\textbf{Hinge Loss}}: A loss function used in binary classification that is differentiable and encourages correct classification with a margin.%
\item \href{\url{https://purdue.brightspace.com/d2l/common/assets/pdfjs-d2l-dist/1.0.14-legacy/web/viewer.html?file=\%252Fcontent\%252Fenforced\%252F1095467-wl.202510.CS.24300.001\%252F05L2-ML\%20Basics.pdf\%253Fou\%253D1095467\&lang=en-us\&container=d2l-fileviewer-rendered-pdf\&fullscreen=d2l-fileviewer-rendered-pdf-dialog\&height=667}#page=27}{\textbf{Logistic Loss}}: A loss function used in binary classification that is differentiable and outputs a probability.%
\item \href{\url{https://purdue.brightspace.com/d2l/common/assets/pdfjs-d2l-dist/1.0.14-legacy/web/viewer.html?file=\%252Fcontent\%252Fenforced\%252F1095467-wl.202510.CS.24300.001\%252F05L2-ML\%20Basics.pdf\%253Fou\%253D1095467\&lang=en-us\&container=d2l-fileviewer-rendered-pdf\&fullscreen=d2l-fileviewer-rendered-pdf-dialog\&height=667}#page=28}{\textbf{Gradient Descent}}: An iterative optimization algorithm used to find the minimum of a function by following the negative gradient.%
\item \href{\url{https://purdue.brightspace.com/d2l/common/assets/pdfjs-d2l-dist/1.0.14-legacy/web/viewer.html?file=\%252Fcontent\%252Fenforced\%252F1095467-wl.202510.CS.24300.001\%252F05L2-ML\%20Basics.pdf\%253Fou\%253D1095467\&lang=en-us\&container=d2l-fileviewer-rendered-pdf\&fullscreen=d2l-fileviewer-rendered-pdf-dialog\&height=667}#page=29}{\textbf{Stochastic Gradient Descent}}: A variant of gradient descent that updates the model parameters using a randomly selected subset of the data.%
\item \href{\url{https://purdue.brightspace.com/d2l/common/assets/pdfjs-d2l-dist/1.0.14-legacy/web/viewer.html?file=\%252Fcontent\%252Fenforced\%252F1095467-wl.202510.CS.24300.001\%252F05L2-ML\%20Basics.pdf\%253Fou\%253D1095467\&lang=en-us\&container=d2l-fileviewer-rendered-pdf\&fullscreen=d2l-fileviewer-rendered-pdf-dialog\&height=667}#page=15}{\textbf{Empirical Risk Minimizer}}: A learning principle that aims to find the model that minimizes the average loss on the training data.%
\end{itemize}

%
\section{Problem Types}%
\label{sec:ProblemTypes}%
\subsection{AI Foundations}%
\label{subsec:AIFoundations}%
\begin{itemize}[leftmargin=*]%
\item \href{\url{https://purdue.brightspace.com/d2l/common/assets/pdfjs-d2l-dist/1.0.14-legacy/web/viewer.html?file=\%252Fcontent\%252Fenforced\%252F1095467-wl.202510.CS.24300.001\%252F02L2-Intro.pdf\%253Fou\%253D1095467\&lang=en-us\&container=d2l-fileviewer-rendered-pdf\&fullscreen=d2l-fileviewer-rendered-pdf-dialog\&height=667}#page=5}{\textbf{Understanding the Nature of Intelligence}}: This problem explores the definition and characteristics of intelligence, both human and artificial. It examines various perspectives on what constitutes intelligence and how it can be measured or simulated.%
\item \href{\url{https://purdue.brightspace.com/d2l/common/assets/pdfjs-d2l-dist/1.0.14-legacy/web/viewer.html?file=\%252Fcontent\%252Fenforced\%252F1095467-wl.202510.CS.24300.001\%252F02L2-Intro.pdf\%253Fou\%253D1095467\&lang=en-us\&container=d2l-fileviewer-rendered-pdf\&fullscreen=d2l-fileviewer-rendered-pdf-dialog\&height=667}#page=5}{\textbf{Defining Artificial Intelligence}}: This problem involves formulating a precise definition of artificial intelligence, considering different perspectives and historical contexts. The challenge lies in capturing the essence of AI while accounting for its evolving nature and diverse applications.%
\item \href{\url{https://purdue.brightspace.com/d2l/common/assets/pdfjs-d2l-dist/1.0.14-legacy/web/viewer.html?file=\%252Fcontent\%252Fenforced\%252F1095467-wl.202510.CS.24300.001\%252F02L2-Intro.pdf\%253Fou\%253D1095467\&lang=en-us\&container=d2l-fileviewer-rendered-pdf\&fullscreen=d2l-fileviewer-rendered-pdf-dialog\&height=667}#page=6}{\textbf{Passing the Turing Test}}: This problem focuses on the implications of passing the Turing Test, a benchmark for machine intelligence. It explores whether passing the test truly indicates machine intelligence or merely the ability to mimic human behavior.%
\item \href{\url{https://purdue.brightspace.com/d2l/common/assets/pdfjs-d2l-dist/1.0.14-legacy/web/viewer.html?file=\%252Fcontent\%252Fenforced\%252F1095467-wl.202510.CS.24300.001\%252F02L2-Intro.pdf\%253Fou\%253D1095467\&lang=en-us\&container=d2l-fileviewer-rendered-pdf\&fullscreen=d2l-fileviewer-rendered-pdf-dialog\&height=667}#page=43}{\textbf{Understanding the AI Grand Challenges}}: This problem explores the reasons behind selecting specific problems, such as autonomous driving, as AI Grand Challenges. It examines the factors that make these problems particularly significant and challenging in the field of AI.%
\item \href{\url{https://purdue.brightspace.com/d2l/common/assets/pdfjs-d2l-dist/1.0.14-legacy/web/viewer.html?file=\%252Fcontent\%252Fenforced\%252F1095467-wl.202510.CS.24300.001\%252F04L1-Linear\%20Algebra\%20Primer.pdf\%253Fou\%253D1095467\&lang=en-us\&container=d2l-fileviewer-rendered-pdf\&fullscreen=d2l-fileviewer-rendered-pdf-dialog\&height=667}#page=17}{\textbf{Understanding Perceptrons}}: The perceptron, a fundamental building block of neural networks, is introduced, illustrating its structure and function.%
\end{itemize}

%
\subsection{AI Development}%
\label{subsec:AIDevelopment}%
\begin{itemize}[leftmargin=*]%
\item \href{\url{https://purdue.brightspace.com/d2l/common/assets/pdfjs-d2l-dist/1.0.14-legacy/web/viewer.html?file=\%252Fcontent\%252Fenforced\%252F1095467-wl.202510.CS.24300.001\%252F02L2-Intro.pdf\%253Fou\%253D1095467\&lang=en-us\&container=d2l-fileviewer-rendered-pdf\&fullscreen=d2l-fileviewer-rendered-pdf-dialog\&height=667}#page=11}{\textbf{Overcoming the Limitations of Search}}: This problem addresses the computational challenges associated with exhaustive search in complex games like chess and Go. It explores the need for alternative approaches, such as machine learning, to overcome the limitations of brute-force search.%
\item \href{\url{https://purdue.brightspace.com/d2l/common/assets/pdfjs-d2l-dist/1.0.14-legacy/web/viewer.html?file=\%252Fcontent\%252Fenforced\%252F1095467-wl.202510.CS.24300.001\%252F02L2-Intro.pdf\%253Fou\%253D1095467\&lang=en-us\&container=d2l-fileviewer-rendered-pdf\&fullscreen=d2l-fileviewer-rendered-pdf-dialog\&height=667}#page=31}{\textbf{Bridging the Gap Between Different AI Domains}}: This problem investigates the similarities and differences between seemingly disparate AI domains, such as game playing and speech synthesis. It explores the potential for transferring knowledge and techniques across different AI applications.%
\item \href{\url{https://purdue.brightspace.com/d2l/common/assets/pdfjs-d2l-dist/1.0.14-legacy/web/viewer.html?file=\%252Fcontent\%252Fenforced\%252F1095467-wl.202510.CS.24300.001\%252F02L2-Intro.pdf\%253Fou\%253D1095467\&lang=en-us\&container=d2l-fileviewer-rendered-pdf\&fullscreen=d2l-fileviewer-rendered-pdf-dialog\&height=667}#page=36}{\textbf{Developing Robust and Generalizable Algorithms}}: This problem focuses on the design of algorithms that can adapt to new situations and generalize their knowledge to unseen scenarios. It explores the challenges in creating algorithms that are not only effective in specific tasks but also robust and adaptable.%
\item \href{\url{https://purdue.brightspace.com/d2l/common/assets/pdfjs-d2l-dist/1.0.14-legacy/web/viewer.html?file=\%252Fcontent\%252Fenforced\%252F1095467-wl.202510.CS.24300.001\%252F02L2-Intro.pdf\%253Fou\%253D1095467\&lang=en-us\&container=d2l-fileviewer-rendered-pdf\&fullscreen=d2l-fileviewer-rendered-pdf-dialog\&height=667}#page=47}{\textbf{Designing Effective AI Systems for Complex Tasks}}: This problem investigates the design and implementation of AI systems for complex tasks, such as game playing and autonomous driving. It explores the various components and techniques involved in building such systems.%
\end{itemize}

%
\subsection{AI Applications}%
\label{subsec:AIApplications}%
\begin{itemize}[leftmargin=*]%
\item \href{\url{https://purdue.brightspace.com/d2l/common/assets/pdfjs-d2l-dist/1.0.14-legacy/web/viewer.html?file=\%252Fcontent\%252Fenforced\%252F1095467-wl.202510.CS.24300.001\%252F02L2-Intro.pdf\%253Fou\%253D1095467\&lang=en-us\&container=d2l-fileviewer-rendered-pdf\&fullscreen=d2l-fileviewer-rendered-pdf-dialog\&height=667}#page=15}{\textbf{Addressing the Challenges of Computer Vision}}: This problem investigates the difficulties in developing computer vision systems that can reliably interpret and understand images. It highlights the gap between initial expectations and the ongoing challenges in this field.%
\item \href{\url{https://purdue.brightspace.com/d2l/common/assets/pdfjs-d2l-dist/1.0.14-legacy/web/viewer.html?file=\%252Fcontent\%252Fenforced\%252F1095467-wl.202510.CS.24300.001\%252F02L2-Intro.pdf\%253Fou\%253D1095467\&lang=en-us\&container=d2l-fileviewer-rendered-pdf\&fullscreen=d2l-fileviewer-rendered-pdf-dialog\&height=667}#page=27}{\textbf{Understanding the Differences Between Game Types}}: This problem explores the differences between various game types, such as chess and soccer, in the context of AI. It examines how these differences affect the design and implementation of AI algorithms for playing these games.%
\item \href{\url{https://purdue.brightspace.com/d2l/common/assets/pdfjs-d2l-dist/1.0.14-legacy/web/viewer.html?file=\%252Fcontent\%252Fenforced\%252F1095467-wl.202510.CS.24300.001\%252F05L2-ML\%20Basics.pdf\%253Fou\%253D1095467\&lang=en-us\&container=d2l-fileviewer-rendered-pdf\&fullscreen=d2l-fileviewer-rendered-pdf-dialog\&height=667}#page=6}{\textbf{Linear Regression}}: Finding the line that best fits a set of data points, minimizing the sum of squared distances between the points and the line.%
\item \href{\url{https://purdue.brightspace.com/d2l/common/assets/pdfjs-d2l-dist/1.0.14-legacy/web/viewer.html?file=\%252Fcontent\%252Fenforced\%252F1095467-wl.202510.CS.24300.001\%252F05L2-ML\%20Basics.pdf\%253Fou\%253D1095467\&lang=en-us\&container=d2l-fileviewer-rendered-pdf\&fullscreen=d2l-fileviewer-rendered-pdf-dialog\&height=667}#page=11}{\textbf{Multiple Linear Regression}}: Extending linear regression to include multiple independent variables to predict a dependent variable.%
\item \href{\url{https://purdue.brightspace.com/d2l/common/assets/pdfjs-d2l-dist/1.0.14-legacy/web/viewer.html?file=\%252Fcontent\%252Fenforced\%252F1095467-wl.202510.CS.24300.001\%252F05L2-ML\%20Basics.pdf\%253Fou\%253D1095467\&lang=en-us\&container=d2l-fileviewer-rendered-pdf\&fullscreen=d2l-fileviewer-rendered-pdf-dialog\&height=667}#page=13}{\textbf{Curve Fitting}}: Fitting a curve to data points, potentially using polynomial functions of higher degrees. This introduces the risk of overfitting.%
\item \href{\url{https://purdue.brightspace.com/d2l/common/assets/pdfjs-d2l-dist/1.0.14-legacy/web/viewer.html?file=\%252Fcontent\%252Fenforced\%252F1095467-wl.202510.CS.24300.001\%252F05L2-ML\%20Basics.pdf\%253Fou\%253D1095467\&lang=en-us\&container=d2l-fileviewer-rendered-pdf\&fullscreen=d2l-fileviewer-rendered-pdf-dialog\&height=667}#page=20}{\textbf{Binary Classification}}: Classifying data points into two categories (+1 or -1) using a linear model and a suitable loss function.%
\end{itemize}

%
\subsection{AI Ethics}%
\label{subsec:AIEthics}%
\begin{itemize}[leftmargin=*]%
\item \href{\url{https://purdue.brightspace.com/d2l/common/assets/pdfjs-d2l-dist/1.0.14-legacy/web/viewer.html?file=\%252Fcontent\%252Fenforced\%252F1095467-wl.202510.CS.24300.001\%252F02L2-Intro.pdf\%253Fou\%253D1095467\&lang=en-us\&container=d2l-fileviewer-rendered-pdf\&fullscreen=d2l-fileviewer-rendered-pdf-dialog\&height=667}#page=19}{\textbf{Mitigating AI Safety Risks}}: This problem examines the potential risks associated with AI systems, such as adversarial attacks and biases. It explores methods for improving AI safety and ensuring responsible development and deployment.%
\item \href{\url{https://purdue.brightspace.com/d2l/common/assets/pdfjs-d2l-dist/1.0.14-legacy/web/viewer.html?file=\%252Fcontent\%252Fenforced\%252F1095467-wl.202510.CS.24300.001\%252F02L2-Intro.pdf\%253Fou\%253D1095467\&lang=en-us\&container=d2l-fileviewer-rendered-pdf\&fullscreen=d2l-fileviewer-rendered-pdf-dialog\&height=667}#page=20}{\textbf{Addressing Ethical Concerns in AI}}: This problem focuses on the ethical implications of AI systems, particularly concerning bias and fairness. It explores the need for responsible AI development and deployment that considers societal impact and avoids perpetuating harmful biases.%
\end{itemize}

%
\subsection{Linear Algebra for AI}%
\label{subsec:LinearAlgebraforAI}%
\begin{itemize}[leftmargin=*]%
\item \href{\url{https://purdue.brightspace.com/d2l/common/assets/pdfjs-d2l-dist/1.0.14-legacy/web/viewer.html?file=\%252Fcontent\%252Fenforced\%252F1095467-wl.202510.CS.24300.001\%252F04L1-Linear\%20Algebra\%20Primer.pdf\%253Fou\%253D1095467\&lang=en-us\&container=d2l-fileviewer-rendered-pdf\&fullscreen=d2l-fileviewer-rendered-pdf-dialog\&height=667}#page=5}{\textbf{Representing Linear Transformations with Matrices}}: Matrices are used to represent linear transformations, which map vectors from one space to another while preserving linear combinations. This is fundamental to many AI algorithms.%
\item \href{\url{https://purdue.brightspace.com/d2l/common/assets/pdfjs-d2l-dist/1.0.14-legacy/web/viewer.html?file=\%252Fcontent\%252Fenforced\%252F1095467-wl.202510.CS.24300.001\%252F04L1-Linear\%20Algebra\%20Primer.pdf\%253Fou\%253D1095467\&lang=en-us\&container=d2l-fileviewer-rendered-pdf\&fullscreen=d2l-fileviewer-rendered-pdf-dialog\&height=667}#page=9}{\textbf{Matrix Arithmetic}}: The slide covers basic matrix operations such as addition, subtraction, scalar multiplication, and multiplication, which are essential for many AI computations.%
\item \href{\url{https://purdue.brightspace.com/d2l/common/assets/pdfjs-d2l-dist/1.0.14-legacy/web/viewer.html?file=\%252Fcontent\%252Fenforced\%252F1095467-wl.202510.CS.24300.001\%252F04L1-Linear\%20Algebra\%20Primer.pdf\%253Fou\%253D1095467\&lang=en-us\&container=d2l-fileviewer-rendered-pdf\&fullscreen=d2l-fileviewer-rendered-pdf-dialog\&height=667}#page=11}{\textbf{Solving Systems of Linear Equations}}: Systems of linear equations can be represented and solved using matrices and their inverses. This is crucial for various AI applications.%
\item \href{\url{https://purdue.brightspace.com/d2l/common/assets/pdfjs-d2l-dist/1.0.14-legacy/web/viewer.html?file=\%252Fcontent\%252Fenforced\%252F1095467-wl.202510.CS.24300.001\%252F04L1-Linear\%20Algebra\%20Primer.pdf\%253Fou\%253D1095467\&lang=en-us\&container=d2l-fileviewer-rendered-pdf\&fullscreen=d2l-fileviewer-rendered-pdf-dialog\&height=667}#page=16}{\textbf{Matrix Transposition}}: Matrix transposition is a fundamental operation that involves swapping rows and columns of a matrix. It's used in various AI algorithms for data manipulation.%
\item \href{\url{https://purdue.brightspace.com/d2l/common/assets/pdfjs-d2l-dist/1.0.14-legacy/web/viewer.html?file=\%252Fcontent\%252Fenforced\%252F1095467-wl.202510.CS.24300.001\%252F04L1-Linear\%20Algebra\%20Primer.pdf\%253Fou\%253D1095467\&lang=en-us\&container=d2l-fileviewer-rendered-pdf\&fullscreen=d2l-fileviewer-rendered-pdf-dialog\&height=667}#page=18}{\textbf{Parallelizing Matrix Multiplication}}: The slide demonstrates how matrix multiplication can be parallelized using domain decomposition, improving computational efficiency.%
\end{itemize}

%
\subsection{Calculus for AI}%
\label{subsec:CalculusforAI}%
\begin{itemize}[leftmargin=*]%
\item \href{\url{https://purdue.brightspace.com/d2l/common/assets/pdfjs-d2l-dist/1.0.14-legacy/web/viewer.html?file=\%252Fcontent\%252Fenforced\%252F1095467-wl.202510.CS.24300.001\%252F04L1-Linear\%20Algebra\%20Primer.pdf\%253Fou\%253D1095467\&lang=en-us\&container=d2l-fileviewer-rendered-pdf\&fullscreen=d2l-fileviewer-rendered-pdf-dialog\&height=667}#page=19}{\textbf{Calculus in AI}}: The slide briefly introduces the importance of calculus, specifically derivatives, partial derivatives, and finding maxima/minima/gradients, in AI algorithms.%
\item \href{\url{https://purdue.brightspace.com/d2l/common/assets/pdfjs-d2l-dist/1.0.14-legacy/web/viewer.html?file=\%252Fcontent\%252Fenforced\%252F1095467-wl.202510.CS.24300.001\%252F04L1-Linear\%20Algebra\%20Primer.pdf\%253Fou\%253D1095467\&lang=en-us\&container=d2l-fileviewer-rendered-pdf\&fullscreen=d2l-fileviewer-rendered-pdf-dialog\&height=667}#page=21}{\textbf{Derivative of the Sigmoid Function}}: The slide shows the derivation of the derivative of the sigmoid function, a crucial component in many AI algorithms, particularly in neural networks.%
\end{itemize}

%
\section{Algorithm Solutions}%
\label{sec:AlgorithmSolutions}%
\subsection{Matrix Operations}%
\label{subsec:MatrixOperations}%
\begin{itemize}[leftmargin=*]%
\item \href{\url{https://purdue.brightspace.com/d2l/common/assets/pdfjs-d2l-dist/1.0.14-legacy/web/viewer.html?file=\%252Fcontent\%252Fenforced\%252F1095467-wl.202510.CS.24300.001\%252F04L1-Linear\%20Algebra\%20Primer.pdf\%253Fou\%253D1095467\&lang=en-us\&container=d2l-fileviewer-rendered-pdf\&fullscreen=d2l-fileviewer-rendered-pdf-dialog\&height=667}#page=9}{\textbf{Matrix Addition}}: $A + B = C$, where $c_{ij} = a_{ij} + b_{ij}$%
\item \href{\url{https://purdue.brightspace.com/d2l/common/assets/pdfjs-d2l-dist/1.0.14-legacy/web/viewer.html?file=\%252Fcontent\%252Fenforced\%252F1095467-wl.202510.CS.24300.001\%252F04L1-Linear\%20Algebra\%20Primer.pdf\%253Fou\%253D1095467\&lang=en-us\&container=d2l-fileviewer-rendered-pdf\&fullscreen=d2l-fileviewer-rendered-pdf-dialog\&height=667}#page=9}{\textbf{Matrix Scalar Multiplication}}: $k \cdot A = B$, where $b_{ij} = k \cdot a_{ij}$%
\item \href{\url{https://purdue.brightspace.com/d2l/common/assets/pdfjs-d2l-dist/1.0.14-legacy/web/viewer.html?file=\%252Fcontent\%252Fenforced\%252F1095467-wl.202510.CS.24300.001\%252F04L1-Linear\%20Algebra\%20Primer.pdf\%253Fou\%253D1095467\&lang=en-us\&container=d2l-fileviewer-rendered-pdf\&fullscreen=d2l-fileviewer-rendered-pdf-dialog\&height=667}#page=14}{\textbf{Matrix Multiplication}}: $C_{ij} = \sum_{k=1}^{n} a_{ik}b_{kj}$%
\item \href{\url{https://purdue.brightspace.com/d2l/common/assets/pdfjs-d2l-dist/1.0.14-legacy/web/viewer.html?file=\%252Fcontent\%252Fenforced\%252F1095467-wl.202510.CS.24300.001\%252F04L1-Linear\%20Algebra\%20Primer.pdf\%253Fou\%253D1095467\&lang=en-us\&container=d2l-fileviewer-rendered-pdf\&fullscreen=d2l-fileviewer-rendered-pdf-dialog\&height=667}#page=10}{\textbf{Matrix Inverse}}: $AA^{-1} = A^{-1}A = I$%
\item \href{\url{https://purdue.brightspace.com/d2l/common/assets/pdfjs-d2l-dist/1.0.14-legacy/web/viewer.html?file=\%252Fcontent\%252Fenforced\%252F1095467-wl.202510.CS.24300.001\%252F04L1-Linear\%20Algebra\%20Primer.pdf\%253Fou\%253D1095467\&lang=en-us\&container=d2l-fileviewer-rendered-pdf\&fullscreen=d2l-fileviewer-rendered-pdf-dialog\&height=667}#page=12}{\textbf{Solving Linear Equations using Matrix Inverses}}: $Ax = b \implies x = A^{-1}b$%
\end{itemize}

%
\subsection{Calculus for Machine Learning}%
\label{subsec:CalculusforMachineLearning}%
\begin{itemize}[leftmargin=*]%
\item \href{\url{https://purdue.brightspace.com/d2l/common/assets/pdfjs-d2l-dist/1.0.14-legacy/web/viewer.html?file=\%252Fcontent\%252Fenforced\%252F1095467-wl.202510.CS.24300.001\%252F04L1-Linear\%20Algebra\%20Primer.pdf\%253Fou\%253D1095467\&lang=en-us\&container=d2l-fileviewer-rendered-pdf\&fullscreen=d2l-fileviewer-rendered-pdf-dialog\&height=667}#page=21}{\textbf{Derivative of Sigmoid Function}}: $\frac{d}{dx} \sigma(x) = \frac{e^{-x}}{(1 + e^{-x})^2}$%
\end{itemize}

%
\subsection{Search Algorithms}%
\label{subsec:SearchAlgorithms}%
\begin{itemize}[leftmargin=*]%
\item \href{\url{https://purdue.brightspace.com/d2l/common/assets/pdfjs-d2l-dist/1.0.14-legacy/web/viewer.html?file=\%252Fcontent\%252Fenforced\%252F1095467-wl.202510.CS.24300.001\%252F02L2-Intro.pdf\%253Fou\%253D1095467\&lang=en-us\&container=d2l-fileviewer-rendered-pdf\&fullscreen=d2l-fileviewer-rendered-pdf-dialog\&height=667}#page=22}{\textbf{Minimax Search}}: A search algorithm for two-player zero-sum games that aims to find the best move for the current player, assuming the opponent also plays optimally. It uses an evaluation function to assign scores to game states and recursively explores the game tree, alternating between maximizing the score for the current player and minimizing the score for the opponent.%
\item \href{\url{https://purdue.brightspace.com/d2l/common/assets/pdfjs-d2l-dist/1.0.14-legacy/web/viewer.html?file=\%252Fcontent\%252Fenforced\%252F1095467-wl.202510.CS.24300.001\%252F02L2-Intro.pdf\%253Fou\%253D1095467\&lang=en-us\&container=d2l-fileviewer-rendered-pdf\&fullscreen=d2l-fileviewer-rendered-pdf-dialog\&height=667}#page=23}{\textbf{Alpha-Beta Pruning}}: A technique to reduce the search space in Minimax search by avoiding the evaluation of branches that cannot affect the optimal decision. It uses alpha (best value for the maximizing player) and beta (best value for the minimizing player) values to prune branches.%
\end{itemize}

%
\subsection{Machine Learning Algorithms}%
\label{subsec:MachineLearningAlgorithms}%
\begin{itemize}[leftmargin=*]%
\item \href{\url{https://purdue.brightspace.com/d2l/common/assets/pdfjs-d2l-dist/1.0.14-legacy/web/viewer.html?file=\%252Fcontent\%252Fenforced\%252F1095467-wl.202510.CS.24300.001\%252F05L2-ML\%20Basics.pdf\%253Fou\%253D1095467\&lang=en-us\&container=d2l-fileviewer-rendered-pdf\&fullscreen=d2l-fileviewer-rendered-pdf-dialog\&height=667}#page=11}{\textbf{Multiple Linear Regression}}: Find the optimal coefficients $\beta^$ that minimize the sum of squared errors between the observed and predicted values. This is given by the formula $\beta^ = (X^TX)^{-1}X^Ty$, where $y$ is the vector of observed values, $X$ is the design matrix, and $\beta$ is the vector of coefficients.%
\end{itemize}

%
\end{document}